\documentclass[aspectratio=169,11pt]{beamer}
\usetheme{Madrid}
\usecolortheme{default}
\setbeamertemplate{navigation symbols}{}
\setbeamertemplate{footline}[frame number]
\usepackage{enumitem}

\usepackage{amsmath,amssymb}
\usepackage{booktabs}
\usepackage{array}

\title{Databricks ML Associate}
\subtitle{Questions 101--200 (Second 100)}
\author{Practice Exam Deck}
\date{}

\begin{document}
	
	\begin{frame}
		\titlepage
	\end{frame}
	
	\begin{frame}{How to Use This Deck}
		\begin{itemize}
			\item This deck contains \textbf{100 new questions} (Q101--Q200), all different from the first 100.
			\item Each question is on its own slide with four options (A--D).
			\item The \textbf{next slide} gives the correct answer and explains why the right option is right and why each wrong option is wrong.
			\item Topics are spread across Feature Store, MLflow, AutoML, data processing, Spark ML, tuning, evaluation, and deployment.
		\end{itemize}
	\end{frame}
	
	% ---------------------------------------------------------------
	\begin{frame}{Question 101}
		Which method retrieves a feature table object by name in the Unity Catalog Feature Store client?
		\begin{enumerate}[label=\Alph*.]
			\item \texttt{fe.get\_table()}
			\item \texttt{fe.load\_table()}
			\item \texttt{fe.read\_table()}
			\item \texttt{fe.fetch\_table()}
		\end{enumerate}
	\end{frame}
	\begin{frame}{Answer 101}
		\textbf{Correct: A} --- \texttt{get\_table()} returns the feature table object.\\
		\vspace{0.5em}
		\textbf{Wrong:}
		\begin{itemize}
			\item B is not an API.
			\item C is not an API.
			\item D is not an API.
		\end{itemize}
	\end{frame}
	
	% ---------------------------------------------------------------
	\begin{frame}{Question 102}
		What is the purpose of a primary key in a Feature Store table?
		\begin{enumerate}[label=\Alph*.]
			\item It uniquely identifies each row for point-in-time lookups
			\item It stores the label
			\item It defines the model type
			\item It sets the cluster size
		\end{enumerate}
	\end{frame}
	\begin{frame}{Answer 102}
		\textbf{Correct: A} --- primary keys identify rows for feature lookups.\\
		\vspace{0.5em}
		\textbf{Wrong:}
		\begin{itemize}
			\item B labels are separate.
			\item C the model type is unrelated.
			\item D cluster size is compute.
		\end{itemize}
	\end{frame}
	
	% ---------------------------------------------------------------
	\begin{frame}{Question 103}
		Which feature table property stores the timestamp used for point-in-time joins?
		\begin{enumerate}[label=\Alph*.]
			\item Timestamp key
			\item Primary key
			\item Partition key
			\item Foreign key
		\end{enumerate}
	\end{frame}
	\begin{frame}{Answer 103}
		\textbf{Correct: A} --- the timestamp key enables point-in-time lookups.\\
		\vspace{0.5em}
		\textbf{Wrong:}
		\begin{itemize}
			\item B is the row identifier.
			\item C is storage layout.
			\item D is relational, not Feature Store.
		\end{itemize}
	\end{frame}
	
	% ---------------------------------------------------------------
	\begin{frame}{Question 104}
		Which method scores a batch DataFrame with a feature-aware model?
		\begin{enumerate}[label=\Alph*.]
			\item \texttt{fe.score\_batch()}
			\item \texttt{fe.predict()}
			\item \texttt{fe.transform()}
			\item \texttt{fe.infer()}
		\end{enumerate}
	\end{frame}
	\begin{frame}{Answer 104}
		\textbf{Correct: A} --- \texttt{score\_batch()} performs feature-aware batch inference.\\
		\vspace{0.5em}
		\textbf{Wrong:}
		\begin{itemize}
			\item B, C, D are not the Feature Store scoring API.
		\end{itemize}
	\end{frame}
	
	% ---------------------------------------------------------------
	\begin{frame}{Question 105}
		What does a FeatureLookup specify?
		\begin{enumerate}[label=\Alph*.]
			\item Which features to join from a feature table
			\item Which model to train
			\item Which cluster to use
			\item Which metric to optimize
		\end{enumerate}
	\end{frame}
	\begin{frame}{Answer 105}
		\textbf{Correct: A} --- FeatureLookup defines the features to retrieve.\\
		\vspace{0.5em}
		\textbf{Wrong:}
		\begin{itemize}
			\item B is model training.
			\item C is compute.
			\item D is tuning.
		\end{itemize}
	\end{frame}
	
	% ---------------------------------------------------------------
	\begin{frame}{Question 106}
		Why should feature tables be created before training sets?
		\begin{enumerate}[label=\Alph*.]
			\item Training sets join features from feature tables
			\item Training sets create feature tables
			\item They are independent
			\item Training sets replace feature tables
		\end{enumerate}
	\end{frame}
	\begin{frame}{Answer 106}
		\textbf{Correct: A} --- training sets depend on feature tables.\\
		\vspace{0.5em}
		\textbf{Wrong:}
		\begin{itemize}
			\item B is reversed.
			\item C is false.
			\item D is false.
		\end{itemize}
	\end{frame}
	
	% ---------------------------------------------------------------
	\begin{frame}{Question 107}
		Which MLflow component stores run metadata such as parameters and metrics?
		\begin{enumerate}[label=\Alph*.]
			\item The tracking server and backend store
			\item The artifact store only
			\item The model registry only
			\item The notebook
		\end{enumerate}
	\end{frame}
	\begin{frame}{Answer 107}
		\textbf{Correct: A} --- the tracking server and backend store hold run metadata.\\
		\vspace{0.5em}
		\textbf{Wrong:}
		\begin{itemize}
			\item B stores artifacts.
			\item C manages models.
			\item D is source code.
		\end{itemize}
	\end{frame}
	
	% ---------------------------------------------------------------
	\begin{frame}{Question 108}
		Which MLflow API logs a parameter?
		\begin{enumerate}[label=\Alph*.]
			\item \texttt{mlflow.log\_param()}
			\item \texttt{mlflow.log\_metric()}
			\item \texttt{mlflow.log\_artifact()}
			\item \texttt{mlflow.set\_tag()}
		\end{enumerate}
	\end{frame}
	\begin{frame}{Answer 108}
		\textbf{Correct: A} --- \texttt{log\_param()} logs parameters.\\
		\vspace{0.5em}
		\textbf{Wrong:}
		\begin{itemize}
			\item B logs metrics.
			\item C logs artifacts.
			\item D sets tags.
		\end{itemize}
	\end{frame}
	
	% ---------------------------------------------------------------
	\begin{frame}{Question 109}
		Which MLflow API sets a tag on a run?
		\begin{enumerate}[label=\Alph*.]
			\item \texttt{mlflow.set\_tag()}
			\item \texttt{mlflow.log\_tag()}
			\item \texttt{mlflow.add\_tag()}
			\item \texttt{mlflow.tag()}
		\end{enumerate}
	\end{frame}
	\begin{frame}{Answer 109}
		\textbf{Correct: A} --- \texttt{set\_tag()} sets a run tag.\\
		\vspace{0.5em}
		\textbf{Wrong:}
		\begin{itemize}
			\item B, C, D are not the API.
		\end{itemize}
	\end{frame}
	
	% ---------------------------------------------------------------
	\begin{frame}{Question 110}
		What does \texttt{mlflow.end\_run()} do?
		\begin{enumerate}[label=\Alph*.]
			\item Ends the active run and records its final status
			\item Deletes the run
			\item Starts a new run
			\item Registers a model
		\end{enumerate}
	\end{frame}
	\begin{frame}{Answer 110}
		\textbf{Correct: A} --- it closes the run and records status.\\
		\vspace{0.5em}
		\textbf{Wrong:}
		\begin{itemize}
			\item B deletion is separate.
			\item C is \texttt{start\_run()}.
			\item D is \texttt{register\_model()}.
		\end{itemize}
	\end{frame}
	
	% ---------------------------------------------------------------
	\begin{frame}{Question 111}
		Which MLflow function loads a model for inference from a registry URI?
		\begin{enumerate}[label=\Alph*.]
			\item \texttt{mlflow.pyfunc.load\_model()}
			\item \texttt{mlflow.load\_run()}
			\item \texttt{mlflow.get\_model()}
			\item \texttt{mlflow.fetch\_model()}
		\end{enumerate}
	\end{frame}
	\begin{frame}{Answer 111}
		\textbf{Correct: A} --- \texttt{pyfunc.load\_model()} loads a model.\\
		\vspace{0.5em}
		\textbf{Wrong:}
		\begin{itemize}
			\item B, C, D are not the loading API.
		\end{itemize}
	\end{frame}
	
	% ---------------------------------------------------------------
	\begin{frame}{Question 112}
		What does an MLflow model signature describe?
		\begin{enumerate}[label=\Alph*.]
			\item The input and output schema of the model
			\item The training algorithm
			\item The cluster size
			\item The experiment name
		\end{enumerate}
	\end{frame}
	\begin{frame}{Answer 112}
		\textbf{Correct: A} --- signatures define input/output schema.\\
		\vspace{0.5em}
		\textbf{Wrong:}
		\begin{itemize}
			\item B is algorithm choice.
			\item C is compute.
			\item D is experiment metadata.
		\end{itemize}
	\end{frame}
	
	% ---------------------------------------------------------------
	\begin{frame}{Question 113}
		Which method registers a model version in Unity Catalog from a run?
		\begin{enumerate}[label=\Alph*.]
			\item \texttt{mlflow.register\_model()}
			\item \texttt{mlflow.log\_model()}
			\item \texttt{mlflow.save\_model()}
			\item \texttt{mlflow.create\_model()}
		\end{enumerate}
	\end{frame}
	\begin{frame}{Answer 113}
		\textbf{Correct: A} --- \texttt{register\_model()} registers a version.\\
		\vspace{0.5em}
		\textbf{Wrong:}
		\begin{itemize}
			\item B logs the artifact.
			\item C saves locally.
			\item D is not the registry API.
		\end{itemize}
	\end{frame}
	
	% ---------------------------------------------------------------
	\begin{frame}{Question 114}
		What is the role of a model alias such as \texttt{@champion}?
		\begin{enumerate}[label=\Alph*.]
			\item It points to a specific model version for deployment
			\item It deletes old versions
			\item It trains the model
			\item It stores metrics
		\end{enumerate}
	\end{frame}
	\begin{frame}{Answer 114}
		\textbf{Correct: A} --- aliases point to versions for serving.\\
		\vspace{0.5em}
		\textbf{Wrong:}
		\begin{itemize}
			\item B aliases do not delete.
			\item C aliases do not train.
			\item D metrics are logged separately.
		\end{itemize}
	\end{frame}
	
	% ---------------------------------------------------------------
	\begin{frame}{Question 115}
		Which MLflow API searches runs in an experiment?
		\begin{enumerate}[label=\Alph*.]
			\item \texttt{mlflow.search\_runs()}
			\item \texttt{mlflow.find\_runs()}
			\item \texttt{mlflow.query\_runs()}
			\item \texttt{mlflow.list\_runs()}
		\end{enumerate}
	\end{frame}
	\begin{frame}{Answer 115}
		\textbf{Correct: A} --- \texttt{search\_runs()} queries runs.\\
		\vspace{0.5em}
		\textbf{Wrong:}
		\begin{itemize}
			\item B, C, D are not the API.
		\end{itemize}
	\end{frame}
	
	% ---------------------------------------------------------------
	\begin{frame}{Question 116}
		Which AutoML API runs regression?
		\begin{enumerate}[label=\Alph*.]
			\item \texttt{databricks.automl.regress()}
			\item \texttt{databricks.automl.classify()}
			\item \texttt{databricks.automl.forecast()}
			\item \texttt{databricks.automl.fit()}
		\end{enumerate}
	\end{frame}
	\begin{frame}{Answer 116}
		\textbf{Correct: A} --- \texttt{regress()} runs regression AutoML.\\
		\vspace{0.5em}
		\textbf{Wrong:}
		\begin{itemize}
			\item B is classification.
			\item C is forecasting.
			\item D is not an API.
		\end{itemize}
	\end{frame}
	
	% ---------------------------------------------------------------
	\begin{frame}{Question 117}
		What does AutoML's leaderboard show?
		\begin{enumerate}[label=\Alph*.]
			\item Candidate models ranked by a primary metric
			\item Cluster costs
			\item Feature tables
			\item Experiment tags
		\end{enumerate}
	\end{frame}
	\begin{frame}{Answer 117}
		\textbf{Correct: A} --- the leaderboard ranks candidate models.\\
		\vspace{0.5em}
		\textbf{Wrong:}
		\begin{itemize}
			\item B is billing.
			\item C is data assets.
			\item D is metadata.
		\end{itemize}
	\end{frame}
	
	% ---------------------------------------------------------------
	\begin{frame}{Question 118}
		Which primary metric does AutoML use for classification by default?
		\begin{enumerate}[label=\Alph*.]
			\item F1
			\item RMSE
			\item MAE
			\item R-squared
		\end{enumerate}
	\end{frame}
	\begin{frame}{Answer 118}
		\textbf{Correct: A} --- classification AutoML optimizes F1.\\
		\vspace{0.5em}
		\textbf{Wrong:}
		\begin{itemize}
			\item B, C, D are regression metrics.
		\end{itemize}
	\end{frame}
	
	% ---------------------------------------------------------------
	\begin{frame}{Question 119}
		Which primary metric does AutoML use for regression by default?
		\begin{enumerate}[label=\Alph*.]
			\item RMSE
			\item F1
			\item Precision
			\item Recall
		\end{enumerate}
	\end{frame}
	\begin{frame}{Answer 119}
		\textbf{Correct: A} --- regression AutoML optimizes RMSE.\\
		\vspace{0.5em}
		\textbf{Wrong:}
		\begin{itemize}
			\item B, C, D are classification metrics.
		\end{itemize}
	\end{frame}
	
	% ---------------------------------------------------------------
	\begin{frame}{Question 120}
		Which file in an MLflow model directory defines the model's flavor and dependencies?
		\begin{enumerate}[label=\Alph*.]
			\item \texttt{MLmodel}
			\item \texttt{model.pkl}
			\item \texttt{README.md}
			\item \texttt{requirements.txt}
		\end{enumerate}
	\end{frame}
	\begin{frame}{Answer 120}
		\textbf{Correct: A} --- \texttt{MLmodel} defines flavor and dependencies.\\
		\vspace{0.5em}
		\textbf{Wrong:}
		\begin{itemize}
			\item B is a serialized model.
			\item C is documentation.
			\item D is one dependency list.
		\end{itemize}
	\end{frame}
	
	% ---------------------------------------------------------------
	\begin{frame}{Question 121}
		Which Spark operation selects a subset of columns?
		\begin{enumerate}[label=\Alph*.]
			\item \texttt{select()}
			\item \texttt{filter()}
			\item \texttt{groupBy()}
			\item \texttt{join()}
		\end{enumerate}
	\end{frame}
	\begin{frame}{Answer 121}
		\textbf{Correct: A} --- \texttt{select()} projects columns.\\
		\vspace{0.5em}
		\textbf{Wrong:}
		\begin{itemize}
			\item B filters rows.
			\item C aggregates.
			\item D joins tables.
		\end{itemize}
	\end{frame}
	
	% ---------------------------------------------------------------
	\begin{frame}{Question 122}
		Which Spark method aggregates rows by a key?
		\begin{enumerate}[label=\Alph*.]
			\item \texttt{groupBy()}
			\item \texttt{select()}
			\item \texttt{withColumn()}
			\item \texttt{orderBy()}
		\end{enumerate}
	\end{frame}
	\begin{frame}{Answer 122}
		\textbf{Correct: A} --- \texttt{groupBy()} groups rows.\\
		\vspace{0.5em}
		\textbf{Wrong:}
		\begin{itemize}
			\item B projects columns.
			\item C adds a column.
			\item D sorts.
		\end{itemize}
	\end{frame}
	
	% ---------------------------------------------------------------
	\begin{frame}{Question 123}
		Which method adds or replaces a column in a Spark DataFrame?
		\begin{enumerate}[label=\Alph*.]
			\item \texttt{withColumn()}
			\item \texttt{select()}
			\item \texttt{filter()}
			\item \texttt{join()}
		\end{enumerate}
	\end{frame}
	\begin{frame}{Answer 123}
		\textbf{Correct: A} --- \texttt{withColumn()} adds or replaces columns.\\
		\vspace{0.5em}
		\textbf{Wrong:}
		\begin{itemize}
			\item B selects existing columns.
			\item C filters rows.
			\item D joins tables.
		\end{itemize}
	\end{frame}
	
	% ---------------------------------------------------------------
	\begin{frame}{Question 124}
		Which method drops rows with nulls in any column?
		\begin{enumerate}[label=\Alph*.]
			\item \texttt{na.drop()}
			\item \texttt{na.fill()}
			\item \texttt{na.replace()}
			\item \texttt{dropna("all")}
		\end{enumerate}
	\end{frame}
	\begin{frame}{Answer 124}
		\textbf{Correct: A} --- \texttt{na.drop()} drops rows with nulls.\\
		\vspace{0.5em}
		\textbf{Wrong:}
		\begin{itemize}
			\item B fills nulls.
			\item C replaces values.
			\item D is the opposite (drop only all-null rows).
		\end{itemize}
	\end{frame}
	
	% ---------------------------------------------------------------
	\begin{frame}{Question 125}
		Which method fills nulls with a constant value?
		\begin{enumerate}[label=\Alph*.]
			\item \texttt{na.fill()}
			\item \texttt{na.drop()}
			\item \texttt{na.replace()}
			\item \texttt{fillna()}
		\end{enumerate}
	\end{frame}
	\begin{frame}{Answer 125}
		\textbf{Correct: A} --- \texttt{na.fill()} fills nulls.\\
		\vspace{0.5em}
		\textbf{Wrong:}
		\begin{itemize}
			\item B drops nulls.
			\item C replaces values.
			\item D is pandas, not Spark.
		\end{itemize}
	\end{frame}
	
	% ---------------------------------------------------------------
	\begin{frame}{Question 126}
		Why is caching a DataFrame useful in iterative ML workflows?
		\begin{enumerate}[label=\Alph*.]
			\item It avoids recomputing the same transformations repeatedly
			\item It trains the model faster
			\item It improves model accuracy
			\item It reduces the feature count
		\end{enumerate}
	\end{frame}
	\begin{frame}{Answer 126}
		\textbf{Correct: A} --- caching avoids recomputation.\\
		\vspace{0.5em}
		\textbf{Wrong:}
		\begin{itemize}
			\item B does not affect model training speed directly.
			\item C does not change accuracy.
			\item D does not reduce features.
		\end{itemize}
	\end{frame}
	
	% ---------------------------------------------------------------
	\begin{frame}{Question 127}
		What is the purpose of repartitioning a Spark DataFrame?
		\begin{enumerate}[label=\Alph*.]
			\item To change the number of partitions for parallelism
			\item To drop nulls
			\item To encode categoricals
			\item To register a model
		\end{enumerate}
	\end{frame}
	\begin{frame}{Answer 127}
		\textbf{Correct: A} --- repartitioning controls parallelism.\\
		\vspace{0.5em}
		\textbf{Wrong:}
		\begin{itemize}
			\item B is null handling.
			\item C is encoding.
			\item D is MLflow.
		\end{itemize}
	\end{frame}
	
	% ---------------------------------------------------------------
	\begin{frame}{Question 128}
		Which operation reduces the number of partitions without a shuffle when possible?
		\begin{enumerate}[label=\Alph*.]
			\item \texttt{coalesce()}
			\item \texttt{repartition()}
			\item \texttt{groupBy()}
			\item \texttt{join()}
		\end{enumerate}
	\end{frame}
	\begin{frame}{Answer 128}
		\textbf{Correct: A} --- \texttt{coalesce()} reduces partitions with minimal shuffle.\\
		\vspace{0.5em}
		\textbf{Wrong:}
		\begin{itemize}
			\item B shuffles.
			\item C aggregates.
			\item D joins.
		\end{itemize}
	\end{frame}
	
	% ---------------------------------------------------------------
	\begin{frame}{Question 129}
		What does \texttt{StringIndexer} output?
		\begin{enumerate}[label=\Alph*.]
			\item A numeric index column for each category
			\item A one-hot vector
			\item A scaled feature
			\item A feature vector
		\end{enumerate}
	\end{frame}
	\begin{frame}{Answer 129}
		\textbf{Correct: A} --- it outputs numeric indices.\\
		\vspace{0.5em}
		\textbf{Wrong:}
		\begin{itemize}
			\item B is OneHotEncoder.
			\item C is StandardScaler.
			\item D is VectorAssembler.
		\end{itemize}
	\end{frame}
	
	% ---------------------------------------------------------------
	\begin{frame}{Question 130}
		What does \texttt{OneHotEncoder} output?
		\begin{enumerate}[label=\Alph*.]
			\item A sparse binary vector for each category
			\item A numeric index
			\item A scaled feature
			\item A label
		\end{enumerate}
	\end{frame}
	\begin{frame}{Answer 130}
		\textbf{Correct: A} --- one-hot outputs binary vectors.\\
		\vspace{0.5em}
		\textbf{Wrong:}
		\begin{itemize}
			\item B is StringIndexer.
			\item C is scaling.
			\item D is the target.
		\end{itemize}
	\end{frame}
	
	% ---------------------------------------------------------------
	\begin{frame}{Question 131}
		Which Spark ML transformer scales features to zero mean and unit variance?
		\begin{enumerate}[label=\Alph*.]
			\item \texttt{StandardScaler}
			\item \texttt{MinMaxScaler}
			\item \texttt{MaxAbsScaler}
			\item \texttt{Normalizer}
		\end{enumerate}
	\end{frame}
	\begin{frame}{Answer 131}
		\textbf{Correct: A} --- StandardScaler uses z-score scaling.\\
		\vspace{0.5em}
		\textbf{Wrong:}
		\begin{itemize}
			\item B scales to a range.
			\item C scales by max absolute value.
			\item D normalizes rows.
		\end{itemize}
	\end{frame}
	
	% ---------------------------------------------------------------
	\begin{frame}{Question 132}
		Which Spark ML transformer scales features to $[0,1]$?
		\begin{enumerate}[label=\Alph*.]
			\item \texttt{MinMaxScaler}
			\item \texttt{StandardScaler}
			\item \texttt{MaxAbsScaler}
			\item \texttt{Normalizer}
		\end{enumerate}
	\end{frame}
	\begin{frame}{Answer 132}
		\textbf{Correct: A} --- MinMaxScaler scales to $[0,1]$.\\
		\vspace{0.5em}
		\textbf{Wrong:}
		\begin{itemize}
			\item B is z-score.
			\item C is max-abs.
			\item D is row normalization.
		\end{itemize}
	\end{frame}
	
	% ---------------------------------------------------------------
	\begin{frame}{Question 133}
		What is the role of a \texttt{Bucketizer} in Spark ML?
		\begin{enumerate}[label=\Alph*.]
			\item It bins continuous features into buckets
			\item It encodes strings
			\item It scales features
			\item It assembles vectors
		\end{enumerate}
	\end{frame}
	\begin{frame}{Answer 133}
		\textbf{Correct: A} --- \texttt{Bucketizer} discretizes continuous features.\\
		\vspace{0.5em}
		\textbf{Wrong:}
		\begin{itemize}
			\item B is StringIndexer.
			\item C is scaling.
			\item D is VectorAssembler.
		\end{itemize}
	\end{frame}
	
	% ---------------------------------------------------------------
	\begin{frame}{Question 134}
		Which Spark ML evaluator is used for binary classification?
		\begin{enumerate}[label=\Alph*.]
			\item \texttt{BinaryClassificationEvaluator}
			\item \texttt{RegressionEvaluator}
			\item \texttt{MulticlassClassificationEvaluator}
			\item \texttt{ClusteringEvaluator}
		\end{enumerate}
	\end{frame}
	\begin{frame}{Answer 134}
		\textbf{Correct: A} --- it evaluates binary classifiers.\\
		\vspace{0.5em}
		\textbf{Wrong:}
		\begin{itemize}
			\item B is for regression.
			\item C is for multiclass.
			\item D is for clustering.
		\end{itemize}
	\end{frame}
	
	% ---------------------------------------------------------------
	\begin{frame}{Question 135}
		Which metric can \texttt{BinaryClassificationEvaluator} compute?
		\begin{enumerate}[label=\Alph*.]
			\item Area under ROC
			\item RMSE
			\item Silhouette
			\item R-squared
		\end{enumerate}
	\end{frame}
	\begin{frame}{Answer 135}
		\textbf{Correct: A} --- it computes AUC and related metrics.\\
		\vspace{0.5em}
		\textbf{Wrong:}
		\begin{itemize}
			\item B is regression.
			\item C is clustering.
			\item D is regression.
		\end{itemize}
	\end{frame}
	
	% ---------------------------------------------------------------
	\begin{frame}{Question 136}
		What does the \texttt{numFolds} parameter in \texttt{CrossValidator} control?
		\begin{enumerate}[label=\Alph*.]
			\item The number of CV folds
			\item The number of trees
			\item The learning rate
			\item The number of features
		\end{enumerate}
	\end{frame}
	\begin{frame}{Answer 136}
		\textbf{Correct: A} --- \texttt{numFolds} sets the number of folds.\\
		\vspace{0.5em}
		\textbf{Wrong:}
		\begin{itemize}
			\item B is \texttt{numTrees}.
			\item C is \texttt{learningRate}.
			\item D is feature count.
		\end{itemize}
	\end{frame}
	
	% ---------------------------------------------------------------
	\begin{frame}{Question 137}
		What does \texttt{parallelism} in \texttt{CrossValidator} control?
		\begin{enumerate}[label=\Alph*.]
			\item How many models are trained concurrently
			\item The number of folds
			\item The number of features
			\item The number of trees
		\end{enumerate}
	\end{frame}
	\begin{frame}{Answer 137}
		\textbf{Correct: A} --- it controls concurrent training.\\
		\vspace{0.5em}
		\textbf{Wrong:}
		\begin{itemize}
			\item B is \texttt{numFolds}.
			\item C is feature count.
			\item D is \texttt{numTrees}.
		\end{itemize}
	\end{frame}
	
	% ---------------------------------------------------------------
	\begin{frame}{Question 138}
		Why is a \texttt{seed} passed to \texttt{CrossValidator}?
		\begin{enumerate}[label=(\Alph*)]
			\item To make fold splits reproducible
			\item To set the number of folds
			\item To set the learning rate
			\item To speed up training
		\end{enumerate}
	\end{frame}
	\begin{frame}{Answer 138}
		\textbf{Correct: A} --- the seed makes splits reproducible.\\
		\vspace{0.5em}
		\textbf{Wrong:}
		\begin{itemize}
			\item B is \texttt{numFolds}.
			\item C is a model parameter.
			\item D does not affect speed.
		\end{itemize}
	\end{frame}
	
	% ---------------------------------------------------------------
	\begin{frame}{Question 139}
		What is the purpose of \texttt{TrainValidationSplit} in Spark ML?
		\begin{enumerate}[label=\Alph*.]
			\item It performs a single train-validation split for tuning
			\item It performs k-fold CV
			\item It trains a model without tuning
			\item It deploys a model
		\end{enumerate}
	\end{frame}
	\begin{frame}{Answer 139}
		\textbf{Correct: A} --- it does a single split for tuning.\\
		\vspace{0.5em}
		\textbf{Wrong:}
		\begin{itemize}
			\item B is \texttt{CrossValidator}.
			\item C is a plain fit.
			\item D is serving.
		\end{itemize}
	\end{frame}
	
	% ---------------------------------------------------------------
	\begin{frame}{Question 140}
		What does a \texttt{ParamGridBuilder} do?
		\begin{enumerate}[label=\Alph*.]
			\item Builds a grid of hyperparameter combinations
			\item Builds a feature vector
			\item Builds a pipeline
			\item Builds a model
		\end{enumerate}
	\end{frame}
	\begin{frame}{Answer 140}
		\textbf{Correct: A} --- it builds a hyperparameter grid.\\
		\vspace{0.5em}
		\textbf{Wrong:}
		\begin{itemize}
			\item B is VectorAssembler.
			\item C is Pipeline.
			\item D is the estimator.
		\end{itemize}
	\end{frame}
	
	% ---------------------------------------------------------------
	\begin{frame}{Question 141}
		Which Spark ML classifier supports multiclass classification natively?
		\begin{enumerate}[label=\Alph*.]
			\item \texttt{LogisticRegression}
			\item \texttt{LinearRegression}
			\item \texttt{KMeans}
			\item \texttt{ALS}
		\end{enumerate}
	\end{frame}
	\begin{frame}{Answer 141}
		\textbf{Correct: A} --- Spark ML logistic regression supports multiclass.\\
		\vspace{0.5em}
		\textbf{Wrong:}
		\begin{itemize}
			\item B is regression.
			\item C is clustering.
			\item D is recommendation.
		\end{itemize}
	\end{frame}
	
	% ---------------------------------------------------------------
	\begin{frame}{Question 142}
		Which Spark ML algorithm is used for collaborative filtering?
		\begin{enumerate}[label=\Alph*.]
			\item \texttt{ALS}
			\item \texttt{KMeans}
			\item \texttt{RandomForest}
			\item \texttt{LogisticRegression}
		\end{enumerate}
	\end{frame}
	\begin{frame}{Answer 142}
		\textbf{Correct: A} --- ALS is used for collaborative filtering.\\
		\vspace{0.5em}
		\textbf{Wrong:}
		\begin{itemize}
			\item B is clustering.
			\item C is classification/regression.
			\item D is classification.
		\end{itemize}
	\end{frame}
	
	% ---------------------------------------------------------------
	\begin{frame}{Question 143}
		Which Spark ML algorithm is a clustering method?
		\begin{enumerate}[label=\Alph*.]
			\item \texttt{KMeans}
			\item \texttt{LinearRegression}
			\item \texttt{RandomForestClassifier}
			\item \texttt{ALS}
		\end{enumerate}
	\end{frame}
	\begin{frame}{Answer 143}
		\textbf{Correct: A} --- KMeans is clustering.\\
		\vspace{0.5em}
		\textbf{Wrong:}
		\begin{itemize}
			\item B is regression.
			\item C is classification.
			\item D is recommendation.
		\end{itemize}
	\end{frame}
	
	% ---------------------------------------------------------------
	\begin{frame}{Question 144}
		What is the purpose of \texttt{maxIter} in Spark ML models?
		\begin{enumerate}[label=\Alph*.]
			\item It caps the number of optimization iterations
			\item It sets the number of trees
			\item It sets the learning rate
			\item It sets the number of features
		\end{enumerate}
	\end{frame}
	\begin{frame}{Answer 144}
		\textbf{Correct: A} --- \texttt{maxIter} caps iterations.\\
		\vspace{0.5em}
		\textbf{Wrong:}
		\begin{itemize}
			\item B is \texttt{numTrees}.
			\item C is \texttt{learningRate}.
			\item D is feature count.
		\end{itemize}
	\end{frame}
	
	% ---------------------------------------------------------------
	\begin{frame}{Question 145}
		Which Spark ML model supports feature importance?
		\begin{enumerate}[label=\Alph*.]
			\item \texttt{RandomForestClassifier}
			\item \texttt{LinearRegression}
			\item \texttt{KMeans}
			\item \texttt{ALS}
		\end{enumerate}
	\end{frame}
	\begin{frame}{Answer 145}
		\textbf{Correct: A} --- tree ensembles expose feature importance.\\
		\vspace{0.5em}
		\textbf{Wrong:}
		\begin{itemize}
			\item B uses coefficients, not importance.
			\item C and D do not expose importance.
		\end{itemize}
	\end{frame}
	
	% ---------------------------------------------------------------
	\begin{frame}{Question 146}
		What does Hyperopt's \texttt{hp.uniform} define?
		\begin{enumerate}[label=\Alph*.]
			\item A uniform search space between two bounds
			\item A log-uniform space
			\item A choice list
			\item A constant
		\end{enumerate}
	\end{frame}
	\begin{frame}{Answer 146}
		\textbf{Correct: A} --- \texttt{hp.uniform} is a uniform range.\\
		\vspace{0.5em}
		\textbf{Wrong:}
		\begin{itemize}
			\item B is \texttt{hp.loguniform}.
			\item C is \texttt{hp.choice}.
			\item D is a constant value.
		\end{itemize}
	\end{frame}
	
	% ---------------------------------------------------------------
	\begin{frame}{Question 147}
		What does Hyperopt's \texttt{hp.choice} define?
		\begin{enumerate}[label=\Alph*.]
			\item A discrete set of options
			\item A uniform range
			\item A log-uniform range
			\item A constant
		\end{enumerate}
	\end{frame}
	\begin{frame}{Answer 147}
		\textbf{Correct: A} --- \texttt{hp.choice} selects from options.\\
		\vspace{0.5em}
		\textbf{Wrong:}
		\begin{itemize}
			\item B is \texttt{hp.uniform}.
			\item C is \texttt{hp.loguniform}.
			\item D is a constant.
		\end{itemize}
	\end{frame}
	
	% ---------------------------------------------------------------
	\begin{frame}{Question 148}
		What does \texttt{SparkTrials} require to run Hyperopt in parallel?
		\begin{enumerate}[label=\Alph*.]
			\item An active Spark cluster
			\item A local Python environment only
			\item A SQL warehouse
			\item A notebook
		\end{enumerate}
	\end{frame}
	\begin{frame}{Answer 148}
		\textbf{Correct: A} --- \texttt{SparkTrials} uses a Spark cluster.\\
		\vspace{0.5em}
		\textbf{Wrong:}
		\begin{itemize}
			\item B is local only.
			\item C is SQL-only.
			\item D is not sufficient.
		\end{itemize}
	\end{frame}
	
	% ---------------------------------------------------------------
	\begin{frame}{Question 149}
		Which Hyperopt search algorithm adapts based on past trials?
		\begin{enumerate}[label=\Alph*.]
			\item TPE
			\item Random search
			\item Grid search
			\item Manual search
		\end{enumerate}
	\end{frame}
	\begin{frame}{Answer 149}
		\textbf{Correct: A} --- TPE adapts using previous results.\\
		\vspace{0.5em}
		\textbf{Wrong:}
		\begin{itemize}
			\item B, C, D do not adapt.
		\end{itemize}
	\end{frame}
	
	% ---------------------------------------------------------------
	\begin{frame}{Question 150}
		Why is Bayesian optimization preferred over grid search in large search spaces?
		\begin{enumerate}[label=\Alph*.]
			\item It uses past results to propose promising configurations
			\item It always finds the global optimum
			\item It requires no objective function
			\item It trains fewer models than grid search by definition
		\end{enumerate}
	\end{frame}
	\begin{frame}{Answer 150}
		\textbf{Correct: A} --- it uses past results to guide the search.\\
		\vspace{0.5em}
		\textbf{Wrong:}
		\begin{itemize}
			\item B no method guarantees the global optimum.
			\item C an objective function is required.
			\item D it can train fewer but not by definition.
		\end{itemize}
	\end{frame}
	
	% ---------------------------------------------------------------
	\begin{frame}{Question 151}
		What is precision?
		\begin{enumerate}[label=\Alph*.]
			\item $TP/(TP+FP)$
			\item $TP/(TP+FN)$
			\item $TN/(TN+FP)$
			\item $(TP+TN)/N$
		\end{enumerate}
	\end{frame}
	\begin{frame}{Answer 151}
		\textbf{Correct: A} --- precision is $TP/(TP+FP)$.\\
		\vspace{0.5em}
		\textbf{Wrong:}
		\begin{itemize}
			\item B is recall.
			\item C is specificity.
			\item D is accuracy.
		\end{itemize}
	\end{frame}
	
	% ---------------------------------------------------------------
	\begin{frame}{Question 152}
		What is recall?
		\begin{enumerate}[label=\Alph*.]
			\item $TP/(TP+FN)$
			\item $TP/(TP+FP)$
			\item $TN/(TN+FP)$
			\item $(TP+TN)/N$
		\end{enumerate}
	\end{frame}
	\begin{frame}{Answer 152}
		\textbf{Correct: A} --- recall is $TP/(TP+FN)$.\\
		\vspace{0.5em}
		\textbf{Wrong:}
		\begin{itemize}
			\item B is precision.
			\item C is specificity.
			\item D is accuracy.
		\end{itemize}
	\end{frame}
	
	% ---------------------------------------------------------------
	\begin{frame}{Question 153}
		What is specificity?
		\begin{enumerate}[label=\Alph*.]
			\item $TN/(TN+FP)$
			\item $TP/(TP+FN)$
			\item $TP/(TP+FP)$
			\item $(TP+TN)/N$
		\end{enumerate}
	\end{frame}
	\begin{frame}{Answer 153}
		\textbf{Correct: A} --- specificity is $TN/(TN+FP)$.\\
		\vspace{0.5em}
		\textbf{Wrong:}
		\begin{itemize}
			\item B is recall.
			\item C is precision.
			\item D is accuracy.
		\end{itemize}
	\end{frame}
	
	% ---------------------------------------------------------------
	\begin{frame}{Question 154}
		What does ROC-AUC measure?
		\begin{enumerate}[label=\Alph*.]
			\item The probability a random positive is ranked above a random negative
			\item The average precision
			\item The accuracy at a fixed threshold
			\item The log loss
		\end{enumerate}
	\end{frame}
	\begin{frame}{Answer 154}
		\textbf{Correct: A} --- ROC-AUC measures ranking probability.\\
		\vspace{0.5em}
		\textbf{Wrong:}
		\begin{itemize}
			\item B is average precision.
			\item C is threshold-specific.
			\item D is log loss.
		\end{itemize}
	\end{frame}
	
	% ---------------------------------------------------------------
	\begin{frame}{Question 155}
		When is PR-AUC preferred over ROC-AUC?
		\begin{enumerate}[label=\Alph*.]
			\item Under strong class imbalance
			\item When classes are balanced
			\item When the target is continuous
			\item When using clustering
		\end{enumerate}
	\end{frame}
	\begin{frame}{Answer 155}
		\textbf{Correct: A} --- PR-AUC is more informative under imbalance.\\
		\vspace{0.5em}
		\textbf{Wrong:}
		\begin{itemize}
			\item B ROC-AUC is fine when balanced.
			\item C is regression.
			\item D is unsupervised.
		\end{itemize}
	\end{frame}
	
	% ---------------------------------------------------------------
	\begin{frame}{Question 156}
		What does a confusion matrix summarize for a binary classifier?
		\begin{enumerate}[label=\Alph*.]
			\item TP, FP, FN, TN counts
			\item Only accuracy
			\item Only precision
			\item The feature importance
		\end{enumerate}
	\end{frame}
	\begin{frame}{Answer 156}
		\textbf{Correct: A} --- it holds the four confusion counts.\\
		\vspace{0.5em}
		\textbf{Wrong:}
		\begin{itemize}
			\item B, C are derived metrics.
			\item D is model interpretation.
		\end{itemize}
	\end{frame}
	
	% ---------------------------------------------------------------
	\begin{frame}{Question 157}
		What is the F1 score?
		\begin{enumerate}[label=\Alph*.]
			\item The harmonic mean of precision and recall
			\item The arithmetic mean of precision and recall
			\item Accuracy under imbalance
			\item The area under the ROC curve
		\end{enumerate}
	\end{frame}
	\begin{frame}{Answer 157}
		\textbf{Correct: A} --- F1 is the harmonic mean of precision and recall.\\
		\vspace{0.5em}
		\textbf{Wrong:}
		\begin{itemize}
			\item B is not F1.
			\item C is not F1.
			\item D is ROC-AUC.
		\end{itemize}
	\end{frame}
	
	% ---------------------------------------------------------------
	\begin{frame}{Question 158}
		Which metric is most sensitive to outliers?
		\begin{enumerate}[label=\Alph*.]
			\item RMSE
			\item MAE
			\item Median absolute error
			\item R-squared
		\end{enumerate}
	\end{frame}
	\begin{frame}{Answer 158}
		\textbf{Correct: A} --- RMSE squares errors, amplifying outliers.\\
		\vspace{0.5em}
		\textbf{Wrong:}
		\begin{itemize}
			\item B is less sensitive.
			\item C is robust.
			\item D is a relative fit measure.
		\end{itemize}
	\end{frame}
	
	% ---------------------------------------------------------------
	\begin{frame}{Question 159}
		What does an $R^2$ of 0 mean?
		\begin{enumerate}[label=\Alph*.]
			\item The model is no better than predicting the mean
			\item The model is perfect
			\item The model is overfit
			\item The model has zero bias
		\end{enumerate}
	\end{frame}
	\begin{frame}{Answer 159}
		\textbf{Correct: A} --- $R^2=0$ matches the mean baseline.\\
		\vspace{0.5em}
		\textbf{Wrong:}
		\begin{itemize}
			\item B is $R^2=1$.
			\item C is not indicated.
			\item D is a bias concept.
		\end{itemize}
	\end{frame}
	
	% ---------------------------------------------------------------
	\begin{frame}{Question 160}
		What does an $R^2$ of 1 mean?
		\begin{enumerate}[label=\Alph*.]
			\item The model explains all variance in the target
			\item The model is no better than the mean
			\item The model has zero variance
			\item The model is underfit
		\end{enumerate}
	\end{frame}
	\begin{frame}{Answer 160}
		\textbf{Correct: A} --- $R^2=1$ is perfect fit.\\
		\vspace{0.5em}
		\textbf{Wrong:}
		\begin{itemize}
			\item B is $R^2=0$.
			\item C is unrelated.
			\item D is not indicated.
		\end{itemize}
	\end{frame}
	
	% ---------------------------------------------------------------
	\begin{frame}{Question 161}
		What is a learning curve?
		\begin{enumerate}[label=\Alph*.]
			\item A plot of train and validation error versus training set size
			\item A plot of accuracy versus threshold
			\item A plot of precision versus recall
			\item A plot of feature importance
		\end{enumerate}
	\end{frame}
	\begin{frame}{Answer 161}
		\textbf{Correct: A} --- learning curves show error vs.\ training size.\\
		\vspace{0.5em}
		\textbf{Wrong:}
		\begin{itemize}
			\item B is a threshold curve.
			\item C is a PR curve.
			\item D is feature importance.
		\end{itemize}
	\end{frame}
	
	% ---------------------------------------------------------------
	\begin{frame}{Question 162}
		What does a persistent gap between training and validation error indicate?
		\begin{enumerate}[label=\Alph*.]
			\item Overfitting / high variance
			\item Underfitting / high bias
			\item Good generalization
			\item Data leakage
		\end{enumerate}
	\end{frame}
	\begin{frame}{Answer 162}
		\textbf{Correct: A} --- a large gap indicates overfitting.\\
		\vspace{0.5em}
		\textbf{Wrong:}
		\begin{itemize}
			\item B underfitting shows high error on both.
			\item C good generalization has a small gap.
			\item D leakage is not shown by a gap alone.
		\end{itemize}
	\end{frame}
	
	% ---------------------------------------------------------------
	\begin{frame}{Question 163}
		What is a validation curve?
		\begin{enumerate}[label=\Alph*.]
			\item A plot of train and validation score versus a hyperparameter
			\item A plot of error versus training size
			\item A plot of precision versus recall
			\item A plot of feature importance
		\end{enumerate}
	\end{frame}
	\begin{frame}{Answer 163}
		\textbf{Correct: A} --- validation curves vary a hyperparameter.\\
		\vspace{0.5em}
		\textbf{Wrong:}
		\begin{itemize}
			\item B is a learning curve.
			\item C is a PR curve.
			\item D is feature importance.
		\end{itemize}
	\end{frame}
	
	% ---------------------------------------------------------------
	\begin{frame}{Question 164}
		What is the bias-variance decomposition of expected test error?
		\begin{enumerate}[label=\Alph*.]
			\item Bias$^2$ + variance + irreducible error
			\item Bias + variance
			\item Training error
			\item Accuracy
		\end{enumerate}
	\end{frame}
	\begin{frame}{Answer 164}
		\textbf{Correct: A} --- the decomposition includes irreducible error.\\
		\vspace{0.5em}
		\textbf{Wrong:}
		\begin{itemize}
			\item B omits irreducible error.
			\item C is training error.
			\item D is a metric.
		\end{itemize}
	\end{frame}
	
	% ---------------------------------------------------------------
	\begin{frame}{Question 165}
		What is the effect of increasing regularization strength?
		\begin{enumerate}[label=\Alph*.]
			\item Bias increases, variance decreases
			\item Bias decreases, variance increases
			\item Both increase
			\item Both decrease
		\end{enumerate}
	\end{frame}
	\begin{frame}{Answer 165}
		\textbf{Correct: A} --- stronger regularization raises bias and lowers variance.\\
		\vspace{0.5em}
		\textbf{Wrong:}
		\begin{itemize}
			\item B is less regularization.
			\item C and D are false.
		\end{itemize}
	\end{frame}
	
	% ---------------------------------------------------------------
	\begin{frame}{Question 166}
		Which method reduces variance by averaging many models?
		\begin{enumerate}[label=\Alph*.]
			\item Bagging
			\item Boosting
			\item Pruning
			\item Early stopping
		\end{enumerate}
	\end{frame}
	\begin{frame}{Answer 166}
		\textbf{Correct: A} --- bagging averages models to reduce variance.\\
		\vspace{0.5em}
		\textbf{Wrong:}
		\begin{itemize}
			\item B boosting reduces bias.
			\item C pruning simplifies a tree.
			\item D stops training early.
		\end{itemize}
	\end{frame}
	
	% ---------------------------------------------------------------
	\begin{frame}{Question 167}
		Which method reduces bias by sequentially correcting errors?
		\begin{enumerate}[label=\Alph*.]
			\item Boosting
			\item Bagging
			\item Pruning
			\item Dropout
		\end{enumerate}
	\end{frame}
	\begin{frame}{Answer 167}
		\textbf{Correct: A} --- boosting corrects errors sequentially.\\
		\vspace{0.5em}
		\textbf{Wrong:}
		\begin{itemize}
			\item B is parallel.
			\item C simplifies trees.
			\item D is a neural net technique.
		\end{itemize}
	\end{frame}
	
	% ---------------------------------------------------------------
	\begin{frame}{Question 168}
		What is the purpose of early stopping?
		\begin{enumerate}[label=\Alph*.]
			\item Stop training when validation performance stops improving
			\item Stop training when training error reaches zero
			\item Stop training after one epoch
			\item Stop training when the model is too simple
		\end{enumerate}
	\end{frame}
	\begin{frame}{Answer 168}
		\textbf{Correct: A} --- early stopping uses validation performance.\\
		\vspace{0.5em}
		\textbf{Wrong:}
		\begin{itemize}
			\item B ignores validation.
			\item C is arbitrary.
			\item D is underfitting.
		\end{itemize}
	\end{frame}
	
	% ---------------------------------------------------------------
	\begin{frame}{Question 169}
		What is a random forest?
		\begin{enumerate}[label=\Alph*.]
			\item An ensemble of decision trees trained on bootstrapped samples with random feature subsets
			\item A single decision tree
			\item A linear model
			\item A clustering algorithm
		\end{enumerate}
	\end{frame}
	\begin{frame}{Answer 169}
		\textbf{Correct: A} --- random forests ensemble trees with bagging.\\
		\vspace{0.5em}
		\textbf{Wrong:}
		\begin{itemize}
			\item B is a single tree.
			\item C is linear regression.
			\item D is unsupervised.
		\end{itemize}
	\end{frame}
	
	% ---------------------------------------------------------------
	\begin{frame}{Question 170}
		What is gradient boosting?
		\begin{enumerate}[label=\Alph*.]
			\item Sequential trees each fitting the previous residual
			\item Parallel trees on bootstrap samples
			\item A clustering method
			\item A dimensionality reduction method
		\end{enumerate}
	\end{frame}
	\begin{frame}{Answer 170}
		\textbf{Correct: A} --- boosting fits residuals sequentially.\\
		\vspace{0.5em}
		\textbf{Wrong:}
		\begin{itemize}
			\item B is bagging.
			\item C is clustering.
			\item D is PCA.
		\end{itemize}
	\end{frame}
	
	% ---------------------------------------------------------------
	\begin{frame}{Question 171}
		What is the main advantage of tree-based models over linear models?
		\begin{enumerate}[label=\Alph*.]
			\item They capture non-linear relationships and interactions
			\item They always train faster
			\item They need no data
			\item They have no hyperparameters
		\end{enumerate}
	\end{frame}
	\begin{frame}{Answer 171}
		\textbf{Correct: A} --- trees capture non-linearity and interactions.\\
		\vspace{0.5em}
		\textbf{Wrong:}
		\begin{itemize}
			\item B is not always true.
			\item C is false.
			\item D trees have many hyperparameters.
		\end{itemize}
	\end{frame}
	
	% ---------------------------------------------------------------
	\begin{frame}{Question 172}
		What is the main advantage of linear models over tree-based models?
		\begin{enumerate}[label=\Alph*.]
			\item Interpretability and extrapolation
			\item They capture all non-linearity
			\item They never overfit
			\item They need no preprocessing
		\end{enumerate}
	\end{frame}
	\begin{frame}{Answer 172}
		\textbf{Correct: A} --- linear models are interpretable and extrapolate.\\
		\vspace{0.5em}
		\textbf{Wrong:}
		\begin{itemize}
			\item B is false; trees capture non-linearity.
			\item C is false.
			\item D is false; scaling is needed.
		\end{itemize}
	\end{frame}
	
	% ---------------------------------------------------------------
	\begin{frame}{Question 173}
		What is k-Nearest Neighbors based on?
		\begin{enumerate}[label=\Alph*.]
			\item Majority/mean of the $k$ nearest training points
			\item A decision boundary learned from trees
			\item A linear combination of features
			\item A neural network
		\end{enumerate}
	\end{frame}
	\begin{frame}{Answer 173}
		\textbf{Correct: A} --- k-NN uses nearest neighbors.\\
		\vspace{0.5em}
		\textbf{Wrong:}
		\begin{itemize}
			\item B is trees.
			\item C is linear models.
			\item D is neural nets.
		\end{itemize}
	\end{frame}
	
	% ---------------------------------------------------------------
	\begin{frame}{Question 174}
		Why must features be scaled before k-NN?
		\begin{enumerate}[label=\Alph*.]
			\item Distance metrics are sensitive to feature scale
			\item k-NN requires categorical features
			\item k-NN cannot use numeric features
			\item Scaling improves accuracy only for trees
		\end{enumerate}
	\end{frame}
	\begin{frame}{Answer 174}
		\textbf{Correct: A} --- k-NN uses distances, so scaling matters.\\
		\vspace{0.5em}
		\textbf{Wrong:}
		\begin{itemize}
			\item B, C are false.
			\item D scaling does not help trees.
		\end{itemize}
	\end{frame}
	
	% ---------------------------------------------------------------
	\begin{frame}{Question 175}
		Which kernel is linear and used for linearly separable SVM problems?
		\begin{enumerate}[label=\Alph*.]
			\item Linear
			\item RBF
			\item Polynomial
			\item Sigmoid
		\end{enumerate}
	\end{frame}
	\begin{frame}{Answer 175}
		\textbf{Correct: A} --- the linear kernel is for linear problems.\\
		\vspace{0.5em}
		\textbf{Wrong:}
		\begin{itemize}
			\item B, C, D are non-linear kernels.
		\end{itemize}
	\end{frame}
	
	% ---------------------------------------------------------------
	\begin{frame}{Question 176}
		What does PCA produce?
		\begin{enumerate}[label=\Alph*.]
			\item Orthogonal components ordered by explained variance
			\item Cluster labels
			\item A decision boundary
			\item A target variable
		\end{enumerate}
	\end{frame}
	\begin{frame}{Answer 176}
		\textbf{Correct: A} --- PCA produces orthogonal components.\\
		\vspace{0.5em}
		\textbf{Wrong:}
		\begin{itemize}
			\item B is clustering.
			\item C is classification.
			\item D is the target.
		\end{itemize}
	\end{frame}
	
	% ---------------------------------------------------------------
	\begin{frame}{Question 177}
		Why is standardization important before PCA?
		\begin{enumerate}[label=\Alph*.]
			\item PCA maximizes variance, so unscaled features dominate
			\item PCA requires categorical features
			\item PCA needs no preprocessing
			\item PCA only works on text
		\end{enumerate}
	\end{frame}
	\begin{frame}{Answer 177}
		\textbf{Correct: A} --- unscaled features can dominate components.\\
		\vspace{0.5em}
		\textbf{Wrong:}
		\begin{itemize}
			\item B, C, D are false.
		\end{itemize}
	\end{frame}
	
	% ---------------------------------------------------------------
	\begin{frame}{Question 178}
		What does the silhouette score measure?
		\begin{enumerate}[label=\Alph*.]
			\item How similar points are to their own cluster vs.\ the nearest other cluster
			\item Model accuracy
			\item Feature importance
			\item The number of clusters
		\end{enumerate}
	\end{frame}
	\begin{frame}{Answer 178}
		\textbf{Correct: A} --- silhouette measures cluster cohesion and separation.\\
		\vspace{0.5em}
		\textbf{Wrong:}
		\begin{itemize}
			\item B is classification.
			\item C is interpretation.
			\item D is a hyperparameter.
		\end{itemize}
	\end{frame}
	
	% ---------------------------------------------------------------
	\begin{frame}{Question 179}
		What range does the silhouette score take?
		\begin{enumerate}[label=\Alph*.]
			\item $[-1, 1]$
			\item $[0, 1]$
			\item $[-\infty, \infty]$
			\item $[0, \infty]$
		\end{enumerate}
	\end{frame}
	\begin{frame}{Answer 179}
		\textbf{Correct: A} --- silhouette ranges from $-1$ to $1$.\\
		\vspace{0.5em}
		\textbf{Wrong:}
		\begin{itemize}
			\item B, C, D are incorrect ranges.
		\end{itemize}
	\end{frame}
	
	% ---------------------------------------------------------------
	\begin{frame}{Question 180}
		How does DBSCAN differ from k-means?
		\begin{enumerate}[label=\Alph*.]
			\item DBSCAN is density-based and finds arbitrary shapes and noise
			\item DBSCAN requires specifying $k$
			\item DBSCAN is for regression
			\item DBSCAN is for classification
		\end{enumerate}
	\end{frame}
	\begin{frame}{Answer 180}
		\textbf{Correct: A} --- DBSCAN is density-based and handles noise.\\
		\vspace{0.5em}
		\textbf{Wrong:}
		\begin{itemize}
			\item B k-means requires $k$.
			\item C, D are supervised tasks.
		\end{itemize}
	\end{frame}
	
	% ---------------------------------------------------------------
	\begin{frame}{Question 181}
		What is anomaly detection?
		\begin{enumerate}[label=\Alph*.]
			\item Identifying rare or abnormal points
			\item Predicting a continuous target
			\item Clustering all points
			\item Reducing dimensionality
		\end{enumerate}
	\end{frame}
	\begin{frame}{Answer 181}
		\textbf{Correct: A} --- anomaly detection finds rare points.\\
		\vspace{0.5em}
		\textbf{Wrong:}
		\begin{itemize}
			\item B is regression.
			\item C is clustering.
			\item D is PCA.
		\end{itemize}
	\end{frame}
	
	% ---------------------------------------------------------------
	\begin{frame}{Question 182}
		Which is an unsupervised anomaly detection method?
		\begin{enumerate}[label=\Alph*.]
			\item Isolation Forest
			\item Logistic Regression
			\item Linear Regression
			\item Decision Tree Classifier
		\end{enumerate}
	\end{frame}
	\begin{frame}{Answer 182}
		\textbf{Correct: A} --- Isolation Forest is unsupervised.\\
		\vspace{0.5em}
		\textbf{Wrong:}
		\begin{itemize}
			\item B, C, D are supervised.
		\end{itemize}
	\end{frame}
	
	% ---------------------------------------------------------------
	\begin{frame}{Question 183}
		What is the purpose of one-hot encoding?
		\begin{enumerate}[label=\Alph*.]
			\item Represent categorical variables as binary columns
			\item Scale numeric features
			\item Impute missing values
			\item Reduce dimensions
		\end{enumerate}
	\end{frame}
	\begin{frame}{Answer 183}
		\textbf{Correct: A} --- one-hot represents categories as binary columns.\\
		\vspace{0.5em}
		\textbf{Wrong:}
		\begin{itemize}
			\item B is scaling.
			\item C is imputation.
			\item D is PCA.
		\end{itemize}
	\end{frame}
	
	% ---------------------------------------------------------------
	\begin{frame}{Question 184}
		Why is label encoding risky for nominal categories in linear models?
		\begin{enumerate}[label=\Alph*.]
			\item It implies an ordinal relationship that does not exist
			\item It creates too many columns
			\item It removes information
			\item It only works for numeric data
		\end{enumerate}
	\end{frame}
	\begin{frame}{Answer 184}
		\textbf{Correct: A} --- label encoding implies false ordering.\\
		\vspace{0.5em}
		\textbf{Wrong:}
		\begin{itemize}
			\item B is one-hot.
			\item C label encoding keeps all categories.
			\item D is false.
		\end{itemize}
	\end{frame}
	
	% ---------------------------------------------------------------
	\begin{frame}{Question 185}
		What is target encoding?
		\begin{enumerate}[label=\Alph*.]
			\item Replacing a category with the mean target value for that category
			\item One-hot encoding
			\item Label encoding
			\item Scaling
		\end{enumerate}
	\end{frame}
	\begin{frame}{Answer 185}
		\textbf{Correct: A} --- target encoding uses category mean targets.\\
		\vspace{0.5em}
		\textbf{Wrong:}
		\begin{itemize}
			\item B, C, D are different techniques.
		\end{itemize}
	\end{frame}
	
	% ---------------------------------------------------------------
	\begin{frame}{Question 186}
		What is a risk of target encoding?
		\begin{enumerate}[label=\Alph*.]
			\item Leakage if not computed with cross-validation
			\item It cannot handle categories
			\item It creates sparse matrices
			\item It requires no target
		\end{enumerate}
	\end{frame}
	\begin{frame}{Answer 186}
		\textbf{Correct: A} --- target encoding can leak without CV.\\
		\vspace{0.5em}
		\textbf{Wrong:}
		\begin{itemize}
			\item B it handles categories.
			\item C sparse matrices are from one-hot.
			\item D it uses the target by definition.
		\end{itemize}
	\end{frame}
	
	% ---------------------------------------------------------------
	\begin{frame}{Question 187}
		What is a Pipeline in Spark ML?
		\begin{enumerate}[label=\Alph*.]
			\item A sequence of stages run as one estimator
			\item A single model
			\item A visualization tool
			\item A storage format
		\end{enumerate}
	\end{frame}
	\begin{frame}{Answer 187}
		\textbf{Correct: A} --- a Pipeline chains stages.\\
		\vspace{0.5em}
		\textbf{Wrong:}
		\begin{itemize}
			\item B is an estimator.
			\item C is visualization.
			\item D is storage.
		\end{itemize}
	\end{frame}
	
	% ---------------------------------------------------------------
	\begin{frame}{Question 188}
		Why should preprocessing be part of a Pipeline?
		\begin{enumerate}[label=\Alph*.]
			\item To fit transforms only on training data and prevent leakage
			\item To speed up training
			\item To reduce features
			\item To register models
		\end{enumerate}
	\end{frame}
	\begin{frame}{Answer 188}
		\textbf{Correct: A} --- pipeline preprocessing prevents leakage.\\
		\vspace{0.5em}
		\textbf{Wrong:}
		\begin{itemize}
			\item B, C, D are not the main reason.
		\end{itemize}
	\end{frame}
	
	% ---------------------------------------------------------------
	\begin{frame}{Question 189}
		Which Spark ML stage assembles multiple columns into a feature vector?
		\begin{enumerate}[label=\Alph*.]
			\item \texttt{VectorAssembler}
			\item \texttt{StringIndexer}
			\item \texttt{OneHotEncoder}
			\item \texttt{StandardScaler}
		\end{enumerate}
	\end{frame}
	\begin{frame}{Answer 189}
		\textbf{Correct: A} --- VectorAssembler builds feature vectors.\\
		\vspace{0.5em}
		\textbf{Wrong:}
		\begin{itemize}
			\item B encodes strings.
			\item C one-hot encodes.
			\item D scales.
		\end{itemize}
	\end{frame}
	
	% ---------------------------------------------------------------
	\begin{frame}{Question 190}
		Which stage converts a string column to numeric indices?
		\begin{enumerate}[label=\Alph*.]
			\item \texttt{StringIndexer}
			\item \texttt{VectorAssembler}
			\item \texttt{OneHotEncoder}
			\item \texttt{StandardScaler}
		\end{enumerate}
	\end{frame}
	\begin{frame}{Answer 190}
		\textbf{Correct: A} --- StringIndexer converts strings to indices.\\
		\vspace{0.5em}
		\textbf{Wrong:}
		\begin{itemize}
			\item B assembles vectors.
			\item C one-hot encodes.
			\item D scales.
		\end{itemize}
	\end{frame}
	
	% ---------------------------------------------------------------
	\begin{frame}{Question 191}
		What is batch inference?
		\begin{enumerate}[label=\Alph*.]
			\item Scoring a large dataset offline
			\item Scoring one record at a time via REST
			\item Training a model
			\item Tuning hyperparameters
		\end{enumerate}
	\end{frame}
	\begin{frame}{Answer 191}
		\textbf{Correct: A} --- batch inference scores offline.\\
		\vspace{0.5em}
		\textbf{Wrong:}
		\begin{itemize}
			\item B is real-time serving.
			\item C is training.
			\item D is tuning.
		\end{itemize}
	\end{frame}
	
	% ---------------------------------------------------------------
	\begin{frame}{Question 192}
		What is real-time inference?
		\begin{enumerate}[label=\Alph*.]
			\item Low-latency prediction via a REST endpoint
			\item Offline scoring
			\item Model training
			\item Data ingestion
		\end{enumerate}
	\end{frame}
	\begin{frame}{Answer 192}
		\textbf{Correct: A} --- real-time inference serves REST predictions.\\
		\vspace{0.5em}
		\textbf{Wrong:}
		\begin{itemize}
			\item B is batch.
			\item C is training.
			\item D is ingestion.
		\end{itemize}
	\end{frame}
	
	% ---------------------------------------------------------------
	\begin{frame}{Question 193}
		What is streaming inference?
		\begin{enumerate}[label=\Alph*.]
			\item Scoring data continuously as it arrives
			\item Scoring a single batch
			\item Training on a stream
			\item Storing predictions
		\end{enumerate}
	\end{frame}
	\begin{frame}{Answer 193}
		\textbf{Correct: A} --- streaming inference scores continuously.\\
		\vspace{0.5em}
		\textbf{Wrong:}
		\begin{itemize}
			\item B is batch.
			\item C is training.
			\item D is storage.
		\end{itemize}
	\end{frame}
	
	% ---------------------------------------------------------------
	\begin{frame}{Question 194}
		Which Databricks feature supports streaming inference?
		\begin{enumerate}[label=\Alph*.]
			\item Delta Live Tables
			\item SQL Warehouse
			\item Single Node cluster
			\item Notebook widgets
		\end{enumerate}
	\end{frame}
	\begin{frame}{Answer 194}
		\textbf{Correct: A} --- Delta Live Tables supports streaming inference.\\
		\vspace{0.5em}
		\textbf{Wrong:}
		\begin{itemize}
			\item B is SQL only.
			\item C is single-node.
			\item D is UI.
		\end{itemize}
	\end{frame}
	
	% ---------------------------------------------------------------
	\begin{frame}{Question 195}
		What is traffic splitting on a serving endpoint?
		\begin{enumerate}[label=\Alph*.]
			\item Routing a percentage of requests to different model versions
			\item Training two models
			\item Storing two models
			\item Deleting old versions
		\end{enumerate}
	\end{frame}
	\begin{frame}{Answer 195}
		\textbf{Correct: A} --- traffic splitting routes requests.\\
		\vspace{0.5em}
		\textbf{Wrong:}
		\begin{itemize}
			\item B is training.
			\item C is storage.
			\item D is deletion.
		\end{itemize}
	\end{frame}
	
	% ---------------------------------------------------------------
	\begin{frame}{Question 196}
		What is model monitoring in production?
		\begin{enumerate}[label=\Alph*.]
			\item Tracking performance and drift over time
			\item Training the model
			\item Registering the model
			\item Deleting the model
		\end{enumerate}
	\end{frame}
	\begin{frame}{Answer 196}
		\textbf{Correct: A} --- monitoring tracks performance and drift.\\
		\vspace{0.5em}
		\textbf{Wrong:}
		\begin{itemize}
			\item B is training.
			\item C is registration.
			\item D is deletion.
		\end{itemize}
	\end{frame}
	
	% ---------------------------------------------------------------
	\begin{frame}{Question 197}
		What is data drift?
		\begin{enumerate}[label=\Alph*.]
			\item A change in the distribution of input features
			\item A change in the feature-target relationship
			\item A model bug
			\item A training error
		\end{enumerate}
	\end{frame}
	\begin{frame}{Answer 197}
		\textbf{Correct: A} --- data drift is an input distribution change.\\
		\vspace{0.5em}
		\textbf{Wrong:}
		\begin{itemize}
			\item B is concept drift.
			\item C is a bug.
			\item D is training error.
		\end{itemize}
	\end{frame}
	
	% ---------------------------------------------------------------
	\begin{frame}{Question 198}
		What is concept drift?
		\begin{enumerate}[label=\Alph*.]
			\item A change in the feature-target relationship
			\item A change in input distribution
			\item A model bug
			\item A data storage issue
		\end{enumerate}
	\end{frame}
	\begin{frame}{Answer 198}
		\textbf{Correct: A} --- concept drift changes the relationship.\\
		\vspace{0.5em}
		\textbf{Wrong:}
		\begin{itemize}
			\item B is data drift.
			\item C is a bug.
			\item D is storage.
		\end{itemize}
	\end{frame}
	
	% ---------------------------------------------------------------
	\begin{frame}{Question 199}
		Which is a fairness concern when deploying ML?
		\begin{enumerate}[label=\Alph*.]
			\item Biased training data leading to discriminatory outcomes
			\item Faster inference
			\item Lower storage cost
			\item Better dashboards
		\end{enumerate}
	\end{frame}
	\begin{frame}{Answer 199}
		\textbf{Correct: A} --- biased data causes discriminatory outcomes.\\
		\vspace{0.5em}
		\textbf{Wrong:}
		\begin{itemize}
			\item B, C, D are performance or cost concerns.
		\end{itemize}
	\end{frame}
	
	% ---------------------------------------------------------------
	\begin{frame}{Question 200}
		Which practice best prevents data leakage in a production pipeline?
		\begin{enumerate}[label=\Alph*.]
			\item Fit preprocessing and the model only on training data inside a Pipeline
			\item Fit preprocessing on the full dataset
			\item Use the test set for feature selection
			\item Scale before splitting
		\end{enumerate}
	\end{frame}
	\begin{frame}{Answer 200}
		\textbf{Correct: A} --- fitting a Pipeline on training data prevents leakage.\\
		\vspace{0.5em}
		\textbf{Wrong:}
		\begin{itemize}
			\item B, C, D cause leakage.
		\end{itemize}
	\end{frame}
	
\end{document}